% Converted from README.md
\subsection{C\# - REST API klijent primer}\label{c---rest-api-klijent-primer}

Aplikacija je primer \href{https://www.redhat.com/en/topics/api/what-is-a-rest-api}{RESTful API} klijenta koji prikazuje trenutnu temperaturu i dnevnu prognozu tako što kontaktira servis \href{https://openweathermap.org/}{OpenWeather}, tačnije njegov \href{https://openweathermap.org/api}{API server} (\emph{Napomena: pokretanje primera zahteva ključ koji aplikacija traži u datoteci \texttt{key.txt}}). Detalji funkcionalnosti ovog servisa nisu od značaja za razumevanje ovog primera. Pojednostavljeno, klijent će serveru poslati HTTP zahtev za odgovarajućim resursom (trenutna temperatura, prognoza, i sl.) i server će poslati objekat sa odgovarajućim informacijama serijalizovanim u JSON. Pogledajmo implementaciju:

\begin{itemize}
\tightlist
\item
  prostor imena \texttt{Common} sadrži klase koje se koriste za deserijalizaciju odgovora servera
\item
  prostor imena \texttt{Services} sadrži klasu \texttt{WeatherService} koja će biti meta naših testova
\item
  glavni prostor imena koristi \texttt{WeatherService} da prikaže trenutnu temperaturu i prognozu za odgovarajući upit
\end{itemize}

Da bismo pisali testove za klasu \texttt{WeatherService}, pogledajmo njen javni interfejs (to što su neke funkcije asinhrone nema uticaj na suštinu primera):

\begin{minted}{csharp}
bool IsDisabled();
async Task<CompleteWeatherData?> GetCurrentDataAsync(string query);
async Task<Forecast?> GetForecastAsync(string query);
\end{minted}

Metod \texttt{IsDisabled} može lako da se testira kreiranjem servisa bez ključa. Druge metode su teže za testiranje jer u sebi rade više poslova: kreiraju HTTP zahtev, šalju zahtev, primaju odgovor i deserijalizuju odgovor. Direktni testovi bi slali zahteve stvarnom OpenWeather servisu, pa bi bili spori i nedeterministički. Takvi testovi ne bi mogli pouzdano da provere ni ponašanje u slučaju nevalidnog ili nekompletnog odgovora servera. Lokalna instanca OpenWeather servera mogla bi da reši deo problema, ali bi testovi i dalje zavisili od dodatnog procesa. Zato ćemo odgovor servera ubrizgati kroz apstrakciju za HTTP komunikaciju. Trenutna implementacija to ne dozvoljava, pa je potrebno da je refaktorišemo.

Pre svega, preimenujmo \texttt{WeatherService} u \texttt{OpenWeatherService} i apstrahujmo interfejs ove klase \texttt{IWeatherService}. Kontaktiranje servera možemo apstrahovati u logiku klase \texttt{WeatherHttpService}. U našim testovima, možemo koristiti implementaciju nalik na onu u klasi \texttt{TestWeatherHttpService}. Ne želimo da \texttt{OpenWeatherService} direktno koristi \texttt{WeatherHttpService} pošto nam to ne omogućava zamenu implementacije klase \texttt{WeatherHttpService} klasom \texttt{TestWeatherHttpService} u našim testovima. Stoga kreirajmo interfejs \texttt{IWeatherHttpService} koji će implementirati klase \texttt{WeatherHttpService} i \texttt{TestWeatherHttpService}. Sada klasa \texttt{OpenWeatherService} može da ima zavisnost na interfejs u konstruktoru (što dodatno omogućava druge povoljnosti kao što je \href{https://www.digitalocean.com/community/tutorials/java-dependency-injection-design-pattern-example-tutorial}{ubrizgavanje zavisnosti}).
